\documentclass[10pt,a4paper]{amsart}
\usepackage{amsmath,amssymb,amsthm,graphicx}
% Retained minimal preamble for the Stacks Project chapter sites-cohomology.tex (Cohomology on Sites).
% Verbatim excerpt of the chapter preamble: Xy-pic setup (native preamble lines 14-18) and the native
% theorem declarations with the mathematical macros (native preamble lines 154-217). Copied mechanically;
% nothing added, removed or redefined. The wrapper class and the amsmath/amssymb/amsthm/graphicx base
% packages are supplied by the assembler, so the class/externaldocument part of the native preamble is omitted.
% For commutative diagrams we use Xy-pic
\usepackage[all]{xy}

% We use 2cell for 2-commutative diagrams.
\xyoption{2cell}
% Theorem environments.
%
\theoremstyle{plain}
\newtheorem{theorem}[subsection]{Theorem}
\newtheorem{proposition}[subsection]{Proposition}
\newtheorem{lemma}[subsection]{Lemma}

\theoremstyle{definition}
\newtheorem{definition}[subsection]{Definition}
\newtheorem{example}[subsection]{Example}
\newtheorem{exercise}[subsection]{Exercise}
\newtheorem{situation}[subsection]{Situation}

\theoremstyle{remark}
\newtheorem{remark}[subsection]{Remark}
\newtheorem{remarks}[subsection]{Remarks}

\numberwithin{equation}{subsection}

% Macros
%
\def\lim{\mathop{\mathrm{lim}}\nolimits}
\def\colim{\mathop{\mathrm{colim}}\nolimits}
\def\Spec{\mathop{\mathrm{Spec}}}
\def\Hom{\mathop{\mathrm{Hom}}\nolimits}
\def\Ext{\mathop{\mathrm{Ext}}\nolimits}
\def\SheafHom{\mathop{\mathcal{H}\!\mathit{om}}\nolimits}
\def\SheafExt{\mathop{\mathcal{E}\!\mathit{xt}}\nolimits}
\def\Sch{\mathit{Sch}}
\def\Mor{\mathop{\mathrm{Mor}}\nolimits}
\def\Ob{\mathop{\mathrm{Ob}}\nolimits}
\def\Sh{\mathop{\mathit{Sh}}\nolimits}
\def\NL{\mathop{N\!L}\nolimits}
\def\CH{\mathop{\mathrm{CH}}\nolimits}
\def\proetale{{pro\text{-}\acute{e}tale}}
\def\etale{{\acute{e}tale}}
\def\QCoh{\mathit{QCoh}}
\def\Ker{\mathop{\mathrm{Ker}}}
\def\Im{\mathop{\mathrm{Im}}}
\def\Coker{\mathop{\mathrm{Coker}}}
\def\Coim{\mathop{\mathrm{Coim}}}

% Boxtimes
%
\DeclareMathSymbol{\boxtimes}{\mathbin}{AMSa}{"02}

%
% Macros for moduli stacks/spaces
%
\def\QCohstack{\mathcal{QC}\!\mathit{oh}}
\def\Cohstack{\mathcal{C}\!\mathit{oh}}
\def\Spacesstack{\mathcal{S}\!\mathit{paces}}
\def\Quotfunctor{\mathrm{Quot}}
\def\Hilbfunctor{\mathrm{Hilb}}
\def\Curvesstack{\mathcal{C}\!\mathit{urves}}
\def\Polarizedstack{\mathcal{P}\!\mathit{olarized}}
\def\Complexesstack{\mathcal{C}\!\mathit{omplexes}}
% \Pic is the operator that assigns to X its picard group, usage \Pic(X)
% \Picardstack_{X/B} denotes the Picard stack of X over B
% \Picardfunctor_{X/B} denotes the Picard functor of X over B
\def\Pic{\mathop{\mathrm{Pic}}\nolimits}
\def\Picardstack{\mathcal{P}\!\mathit{ic}}
\def\Picardfunctor{\mathrm{Pic}}
\def\Deformationcategory{\mathcal{D}\!\mathit{ef}}

% --- appended 2026-10-12 (Chief): the project own definitions of the reference/slogan/history
% display environments, copied mechanically from the shared Stacks preamble lines 8-12
% (\usepackage{verbatim} plus the three \newenvironment lines, which the chapter preamble keeps
% verbatim from the project preamble). Nothing added, removed or redefined: these environments
% carry no mathematical content and the project itself typesets them as comments. Without these
% three lines a counted statement whose own text uses them fails with Environment slogan undefined.
\usepackage{verbatim}
\newenvironment{reference}{\comment}{\endcomment}
\newenvironment{slogan}{\comment}{\endcomment}
\newenvironment{history}{\comment}{\endcomment}
\usepackage{cleveref}
\crefname{originalsection}{original section}{original sections}
\Crefname{originalsection}{Original section}{Original sections}
\makeatletter
\newlabel{section-compare}{{0EWK}{source}{}{sourcepointer-1}{}}
\newlabel{section-compare@cref}{{[originalsection][0][]0EWK}{[1][0][]source}{}{}{}}
\makeatother
\crefname{originallemma}{original lemma}{original lemmas}
\Crefname{originallemma}{Original lemma}{Original lemmas}
\makeatletter
\newlabel{lemma-compare-cohomology-general}{{0EZH}{source}{}{sourcepointer-2}{}}
\newlabel{lemma-compare-cohomology-general@cref}{{[originallemma][0][]0EZH}{[1][0][]source}{}{}{}}
\makeatother
\crefname{originallemma}{original lemma}{original lemmas}
\Crefname{originallemma}{Original lemma}{Original lemmas}
\makeatletter
\newlabel{lemma-cohomological-descent-general}{{0EZG}{source}{}{sourcepointer-3}{}}
\newlabel{lemma-cohomological-descent-general@cref}{{[originallemma][0][]0EZG}{[1][0][]source}{}{}{}}
\makeatother
\crefname{originallemma}{original lemma}{original lemmas}
\Crefname{originallemma}{Original lemma}{Original lemmas}
\makeatletter
\newlabel{lemma-Leray-unbounded}{{0D6H}{source}{}{sourcepointer-4}{}}
\newlabel{lemma-Leray-unbounded@cref}{{[originallemma][0][]0D6H}{[1][0][]source}{}{}{}}
\makeatother
\title{Property (V) implies the comparison of cohomology in degrees up to n}
\author{The Stacks Project Authors}
\date{stacks-project:sites-cohomology.tex:local-snapshot-20261011; GNU FDL 1.2 or later}
\begin{document}
\maketitle
\noindent\textit{Editorial note.} All delivered mathematical text is the author's own, copied from the pinned Stacks Project chapter file sites-cohomology.tex; no statement or proof step was rewritten, reordered or removed.\par\medskip
\noindent\textit{Editorial note.} The proof is the author's complete argument; the parenthetical '(small detail omitted)' concerns the passage from the degeneration of the spectral sequence to the degreewise isomorphism, a routine identification, and the sentence is preserved verbatim in the counted span.\par\medskip
\noindent\textit{Editorial note.} Every labelled reference occurring inside the counted spans is carried in assembly.external\_references as a non-load-bearing navigation pointer (same-chapter ones at their definition site, cross-chapter ones at the invocation inside the counted span), so no load-bearing mathematics is replaced by a stub. No bibliography citation occurs inside the counted spans.\par\medskip
\noindent\textit{Editorial note.} The checked chapter list for same-name collisions is the set of stacks chapters staged in this repository: chow, cohomology, divisors, etale-cohomology, flat, more-algebra, more-morphisms, morphisms, spaces-more-morphisms, varieties together with this card's chapter sites-cohomology.\par\medskip
\noindent\textit{Editorial note.} The preamble is the chapter's own retained minimal preamble (sources/stacks-sites-cohomology/retained-preamble.tex), a verbatim excerpt of the native preamble: the Xy-pic setup and the native theorem declarations with the mathematical macros.\par\medskip
\noindent\textit{Allowed standard prerequisites (editorial).}\begin{itemize}
\item The comparison situation of the section. Same chapter, Situation situation-compare and the standing setup of Section section-compare. Supplies epsilon and the subcategories in which the object of the statement lives.
\item The unbounded Leray spectral sequence. Same chapter, Lemma lemma-Leray-unbounded. Supplies the identification of the cohomology on the coarser site with the cohomology of the derived pushforward.
\item The spectral sequence of the derived pushforward. Same chapter, the spectral sequence with second page given by the higher direct images of the cohomology sheaves. Supplies the spectral sequence whose degeneration is used.
\end{itemize}\medskip
\ifdefined\hypertarget\hypertarget{sourcepointer-1}{}\fi
\noindent\textit{Source reference.} Original section 0EWK: sources/stacks-sites-cohomology/pinned-original-source.tex, source lines 6177--6177. Navigation pointer invoked at proof line 6813; the counted span contains the pointer only and the cited result is not reproduced inside it, so the pointer is registered as a non-load-bearing navigation stub in the established form for Stacks material. No bibliography citation is involved.\par\medskip
\ifdefined\hypertarget\hypertarget{sourcepointer-2}{}\fi
\noindent\textit{Source reference.} Original lemma 0EZH: sources/stacks-sites-cohomology/pinned-original-source.tex, source lines 7425--7425. Navigation pointer invoked at proof line 6905; the counted span contains the pointer only and the cited result is not reproduced inside it, so the pointer is registered as a non-load-bearing navigation stub in the established form for Stacks material. No bibliography citation is involved.\par\medskip
\ifdefined\hypertarget\hypertarget{sourcepointer-3}{}\fi
\noindent\textit{Source reference.} Original lemma 0EZG: sources/stacks-sites-cohomology/pinned-original-source.tex, source lines 7361--7361. Navigation pointer invoked at proof line 6906; the counted span contains the pointer only and the cited result is not reproduced inside it, so the pointer is registered as a non-load-bearing navigation stub in the established form for Stacks material. No bibliography citation is involved.\par\medskip
\ifdefined\hypertarget\hypertarget{sourcepointer-4}{}\fi
\noindent\textit{Source reference.} Original lemma 0D6H: sources/stacks-sites-cohomology/pinned-original-source.tex, source lines 4403--4403. Navigation pointer invoked at proof line 7038; the counted span contains the pointer only and the cited result is not reproduced inside it, so the pointer is registered as a non-load-bearing navigation stub in the established form for Stacks material. No bibliography citation is involved.\par\medskip
\begin{situation}
\label{situation-compare}
With $\mathcal{C}$, $\tau$, and $\tau'$ as in Section \ref{section-compare}.
Assume we are given a subset $\mathcal{P} \subset \text{Arrows}(\mathcal{C})$
and for every object $X$ of $\mathcal{C}$ we are given a weak Serre subcategory
$\mathcal{A}'_X \subset \textit{Ab}(\mathcal{C}_{\tau'}/X)$.
We make the following assumption:
\begin{enumerate}
\item
\label{item-base-change-P}
given $f : X \to Y$ in $\mathcal{P}$ and $Y' \to Y$ general,
then $X \times_Y Y'$ exists and $X \times_Y Y' \to Y'$ is in $\mathcal{P}$,
\item
\label{item-restriction-A}
$f_{\tau'}^{-1}$ sends $\mathcal{A}'_Y$ into $\mathcal{A}'_X$
for any morphism $f : X \to Y$ of $\mathcal{C}$,
\item
\label{item-A-sheaf}
given $X$ in $\mathcal{C}$ and $\mathcal{F}'$ in $\mathcal{A}'_X$, then
$\mathcal{F}'$ satisfies the sheaf condition for $\tau$-coverings, i.e.,
$\mathcal{F}' = \epsilon_{X, *}\epsilon_X^{-1}\mathcal{F}'$,
\item
\label{item-A-and-P}
if $f : X \to Y$ in $\mathcal{P}$ and
$\mathcal{F}' \in \Ob(\mathcal{A}'_X)$, then
$R^if_{\tau', *}\mathcal{F}' \in \Ob(\mathcal{A}'_Y)$
for $i \geq 0$.
\item
\label{item-refine-tau-by-P}
if $\{U_i \to U\}_{i \in I}$ is a $\tau$-covering, then there exist
\begin{enumerate}
\item a $\tau'$-covering $\{V_j \to U\}_{j \in J}$,
\item a $\tau$-covering $\{f_j : W_j \to V_j\}$ consisting
of a single $f_j \in \mathcal{P}$, and
\item a $\tau'$-covering $\{W_{jk} \to W_j\}_{k \in K_j}$
\end{enumerate}
such that $\{W_{jk} \to U\}_{j \in J, k \in K_j}$ is a refinement of
$\{U_i \to U\}_{i \in I}$.
\end{enumerate}
\end{situation}

\noindent
Our next goal is to prove Lemmas \ref{lemma-compare-cohomology-general} and
\ref{lemma-cohomological-descent-general}. We will do this by an induction
argument using the following induction hypothesis.

\medskip\noindent
$(V_n)$ For $X$ in $\mathcal{C}$ and $\mathcal{F}$ in $\mathcal{A}_X$
we have $R^i\epsilon_{X, *}\mathcal{F} = 0$ for $1 \leq i \leq n$.

\begin{lemma}
\label{lemma-V-implies-cohomology-general}
In Situation \ref{situation-compare} if $(V_n)$ holds, then
for $X$ in $\mathcal{C}$ and $L \in D(\mathcal{C}_{\tau'}/X)$
with $H^i(L) = 0$ for $i < 0$ and $H^i(L)$ in $\mathcal{A}'_X$
for $0 \leq i \leq n$ we have
$H^n_{\tau'}(X, L) = H^n_\tau(X, \epsilon_X^{-1}L)$.
\end{lemma}

\begin{proof}
By Lemma \ref{lemma-Leray-unbounded} we have
$H^n_\tau(X, \epsilon_X^{-1}L) =
H^n_{\tau'}(X, R\epsilon_{X, *}\epsilon_X^{-1}L)$.
There is a spectral sequence
$$
E_2^{p, q} = R^p\epsilon_{X, *}\epsilon_X^{-1}H^q(L)
$$
converging to $H^{p + q}(R\epsilon_{X, *}\epsilon_X^{-1}L)$.
By $(V_n)$ we have the vanishing of $E_2^{p, q}$ for
$0 < p \leq n$ and $0 \leq q \leq n$. Thus
$E_2^{0, q} = \epsilon_{X, *}\epsilon_X^{-1}H^q(L) = H^q(L)$
are the only nonzero terms $E_2^{p, q}$ with $p + q \leq n$.
It follows that the map
$$
L \longrightarrow R\epsilon_{X, *}\epsilon_X^{-1}L
$$
is an isomorphism in degrees $\leq n$ (small detail omitted).
Hence we find that
$H^i_{\tau'}(X, L) = H^i_{\tau'}(X, R\epsilon_{X, *}\epsilon_X^{-1}L)$
for $i \leq n$. Thus the lemma is proved.
\end{proof}

\end{document}
